\section{Comparison between the Covariances under a Specific Parametric Form}
\label{app:comparison}

Consider the specification of $\Prob(y = j\,|\,\x, z)$ given by 
\begin{align}
    g_j(\x, z; \thetavec^*) ~=~ \frac{e^{\check{\thetavec}^{\top}\x_{(j)}+ \eta\ind_{(z = j)}}}{1+\sum_{\ell \in \Kscr} e^{\check{\thetavec}^{\top}\x_{(\ell)}+ \eta\ind_{(z = \ell)}}}, \quad \forall j \in \Kscr.
    \label{eqn:mnl-g}
\end{align} 
Let us define the random variable
$
\rvec(\x, z) := \sum_{j\in \Kscr}\big(\sigma_{j}(\x; \betavec^*) 
 - g_{j}(\x, z; \thetavec^*)\big)\x_{(j)}
$. 
The next result follows. 

\begin{prop}[Projection Error Decomposition under MNL]
\label{prop:error-mnl}
    If \eqref{eqn:mnl-g} holds, under Assumptions \ref{assumption:regularity-consistency} and \ref{assumption:regularity-normality},  $\check{\Jvec} -  \Gammavec\Lambdavec\Gammavec^{\top}$ equals to:
  \begin{align}
       &\E\left[\rvec(\x, z)\rvec(\x, z)^{\top}\right]
        ~=~
     \underbrace{\E_{\x}\left[
      \left(\sum_{j\in \Kscr}\big(\sigma_{j}(\x; \betavec^*) 
 - \Prob(y = j\,|\,\x)\big)\x_{(j)}
     \right)\left(\sum_{j\in \Kscr}\big(\sigma_{j}(\x; \betavec^*) 
 - \Prob(y = j\,|\,\x)\big)\x_{(j)}
     \right)^{\top}
     \right]}_{\text{\upshape{The degree of model misspecification}}}\notag\\
     &~~~~~~~~~~~~~~~~~~~~+\underbrace{\E_{\x,z}\left[
      \left(\sum_{j\in \Kscr}\big(g_{j}(\x, z; \thetavec^*)
 - \Prob(y = j\,|\,\x)\big)\x_{(j)}
     \right)\left(\sum_{j\in \Kscr}\big(g_{j}(\x, z; \thetavec^*)
 - \Prob(y = j\,|\,\x)\big)\x_{(j)}
     \right)^{\top}
     \right]}_{\text{\upshape{The covariance of $\sum_{j\in \Kscr}g_{j}(\x, z; \thetavec^*))\x_{(j)}$}}}\succ 0.
     \label{eqn:decomposition-variance}
  \end{align}
\end{prop}

In this case, the difference $ \check{\Jvec} - \Gammavec\Lambdavec\Gammavec^{\top} $ decomposes neatly into two terms: one term reflects the extent to which the MNL choice model is misspecified, and the other term captures the covariance of $ \sum_{j \in \Kscr} g_{j}(\x, z; \thetavec^*) \x_{(j)} $,
because  
$
\E_{z}\left[
g_{j}(\x, z; \thetavec^*) \,\big|\,\x
\right]\,=\,\Prob(y = j\,|\,\x).
$
This decomposition offers valuable insights into the conditions under which AAE achieves variance reduction compared to $ \hat{\betavec}^{\sf P}$. 
Specifically, if model misspecification exists, so that the matrix from above is positive definite, or if the auxiliary label $ z $ is sufficiently informative, such that the associated covariance matrix is positive definite, then $\check{\Jvec} \succ \Gammavec\Lambdavec\Gammavec^{\top}$ holds, establishing AAE’s dominance. 
Furthermore, as model misspecification increases or $ z $ becomes more informative, this dominance is amplified, resulting in a larger variance reduction by AAE.

\begin{figure}[!ht]
    \centering
    \includegraphics[width=120mm]{fig/eigen_prob_plot.pdf}
    \caption{Minimum Eigenvalues and Absolute Probability Differences}
    \label{fig:eigen}
    \begin{minipage}{\textwidth}
    \footnotesize
    \begin{tabular}{p{1cm} p{15cm}}\hspace{-1mm}\textbf{Notes}: & \setlength{\baselineskip}{1.2\baselineskip} The plot shows averages of minimum eigenvalues (blue curve) and absolute probability differences (red curve) across 50 instances when $\eta$ changes.
         The ribbons give the maximum and minimum values for either the minimum eigenvalue or the absolute probability difference corresponding to each $\eta$ value. The vertical axes are adjusted for the minimum eigenvalue or the absolute probability difference corresponding for better visualization. 
    \end{tabular}
  \end{minipage}
\end{figure}

Using this proposition, we can compute $ \check{\Jvec} -  \Gammavec\Lambdavec\Gammavec^{\top} = (1 - \alpha)\alpha(2p - 1)^2 > 0 $ in the context of Example \ref{examp:ai-bias}, given that $ p \neq 1/2 $ and $ \alpha \notin \{0,1\} $. Note that, in this simplified setting, there is no model misspecification as in \eqref{eqn:decomposition-variance}. When $ p = 1/2 $, $z$ becomes pure noise, and when $ \alpha \in \{0,1\} $, $ z $ is constant. In both cases, $ z $ is uninformative. Outside of these cases, we find that AAE performs better than using only the primary data. The advantage becomes more evident if we send $p \rightarrow 0$ or $1$. 
In these cases, $z$ becomes more informative. 

Next, we expand on this discussion, and numerically illustrate $ \check{\Jvec} -  \Gammavec\Lambdavec\Gammavec^{\top} $ under the specification of \eqref{eqn:mnl-g}, 
and plot the results in Figure \ref{fig:eigen}. Specifically, we assume that 
$
\Prob(z = j \mid \x) = {e^{\bm{\zeta}^{\top}\x_{(j)}}}/\big(\sum_{\ell \in \Kscr} e^{\bm{\zeta}^{\top}\x_{(\ell)}}\big),\,\forall j\in \Kscr
$,
with $ d = 5 $ and $ k = 2 $.
Using the specification in \eqref{eqn:mnl-g}, we vary $ \eta $ across the set $ \{0.01, 0.1, 0.5, 1, 2, 3, 4, 5, 6, 7, 8\} $, generating 50 instances for each $ \eta $ value and calculating the minimum eigenvalue of $ \check{\Jvec} -  \Gammavec\Lambdavec\Gammavec^{\top} = \E\left[\rvec(\x, z)\rvec(\x, z)^{\top}\right] $. In each instance, we sample each component of $ \check{\thetavec} $ and $ \bm{\zeta} $ from $ U[-2, 2] $ and each component of $ \x $ from $ U[-1, 1] $, independently of one another.
In Figure \ref{fig:eigen}, we plot the average minimum eigenvalues across the 50 instances for each $ \eta $ (blue curve) alongside the average \textit{absolute probability differences} (red curve). The absolute probability difference here is defined as
\[
\E_{\x, z}\left[\frac{1}{k} \sum_{j \in \Kscr} \left| \sigma_{j}(\x; \betavec^*) - g_{j}(\x, z; \thetavec^*) \right|\right],
\]
which serves as an indicator of the distance between $ \sigma_{j}(\x; \betavec^*) $ and $ g_{j}(\x, z; \thetavec^*) $, thus intuitively reflecting the importance of $ z $.
Note that one can also use the KL-divergence as the metric, but we find its scale difficult to comprehend intuitively compared with the absolute probability difference. 

As $ \eta $ increases, the alignment between the random realizations of $ y $ and $ z $ also increases, signifying the growing value of the AI-generated label $ z $. This effect is illustrated in the red curve. 
The minimum eigenvalue also rises in a similar trend. Notably, the minimum eigenvalue remains consistently positive across all instances, even when $ \eta = 0.01 $, very close to zero. In this case, $z$ is nearly a noise; yet in the worst case across the 50 instances, the minimum eigenvalue is still approximately $ 1 \times 10^{-6} > 0 $, indicating that as long as $ z $ contains a minimal amount of information, we generally expect the second part of Proposition \ref{prop:weak-dom} to hold.

\subsection{Additional Proofs}

\noindent
\textbf{Proof of Proposition \ref{prop:error-mnl}.}
To begin, let us recall that $
\rvec ~=~ \sum_{j\in \Kscr}\big(\sigma_{j}(\x; \betavec^*) 
 - g_{j}(\x, z; \thetavec^*)\big)\x_{(j)}
$, where we suppress the dependency on $\x$ and $z$ in $\rvec$ for notation brevity. 
Recycling the notations used in the proof of Proposition \ref{prop:weak-dom}, we have 
$$\wvec~=~\sum_{j\in \Kscr}\big(\ind_{\{y = j\}} - \sigma_{j}(\x; \betavec^*)\big)\x_{(j)},
$$
and 
$$
\uvec = \nabla_{\thetavec} \log g_y(\x, z, \thetavec^*) = \sum_{j\in \Kscr}\big(\ind_{\{y = j\}} - g_{j}(\x, z; \thetavec^*)\big)\begin{pmatrix}
    \x_{(j)}\\\\\ind_{\{z = j\}}
\end{pmatrix}=
\begin{pmatrix}
    \tilde{\wvec}\\\\u_0
\end{pmatrix},
$$
where $\tilde{\wvec} := \sum_{j\in \Kscr}\big(\ind_{\{y = j\}} - g_{j}(\x, z; \thetavec^*)\big)\x_{(j)}$ and $u_0 :=\sum_{j\in \Kscr}\big(\ind_{\{y = j\}} - g_{j}(\x, z; \thetavec^*)\big)\ind_{\{z = j\}}$.
Therefore, $\rvec = \tilde{\wvec} - \wvec$. 
Also, recall that $$
\check{\Jvec} - \Gammavec\Lambdavec\Gammavec^{\top}~=~
\E\left[\wvec \wvec^{\top} \right] - \E\left[\wvec \uvec^{\top} \right]\E\left[\uvec \uvec^{\top} \right]^{-1}\E\left[\uvec \wvec^{\top}\right].
$$
We examine each term separately. 
We note 
\begin{align*}
 \wvec \uvec^{\top} = \wvec 
 \begin{pmatrix}
    \tilde{\wvec}^{\top} & u_0
 \end{pmatrix}
 = (\tilde{\wvec} - \rvec) 
 \begin{pmatrix}
    \tilde{\wvec}^{\top} & u_0
 \end{pmatrix}
 =  \begin{pmatrix}
    \tilde{\wvec}\tilde{\wvec}^{\top} & u_0\tilde{\wvec}   
 \end{pmatrix} - 
 \begin{pmatrix}
    \rvec\tilde{\wvec}^{\top} & u_0\rvec   
 \end{pmatrix},
\end{align*}
and 
$$
\uvec \uvec^{\top} ~=~
\begin{pmatrix}
    \tilde{\wvec}\tilde{\wvec}^{\top} & u_0\tilde{\wvec}\\
    u_0\tilde{\wvec}^{\top} & u_0^2
\end{pmatrix}.
$$
Thus, 
\begin{align*}
    &\E\left[\wvec \uvec^{\top} \right]\E\left[\uvec \uvec^{\top} \right]^{-1}\E\left[\uvec \wvec^{\top}\right]~=~
    \underbrace{\E\left[\begin{pmatrix}
    \tilde{\wvec}\tilde{\wvec}^{\top} & u_0\tilde{\wvec}   
 \end{pmatrix} \right]
 \E\left[
 \begin{pmatrix}
    \tilde{\wvec}\tilde{\wvec}^{\top} & u_0\tilde{\wvec}\\
    u_0\tilde{\wvec}^{\top} & u_0^2
\end{pmatrix}\right]^{-1}
\E\left[\begin{pmatrix}
    \tilde{\wvec}\tilde{\wvec}^{\top} \\ u_0\tilde{\wvec}^{\top}  
 \end{pmatrix} \right]}_{=:\Bvec_1}\\
 &-\underbrace{\E\left[\begin{pmatrix}
    \rvec\tilde{\wvec}^{\top} & u_0\rvec  
 \end{pmatrix} \right]
 \E\left[
 \begin{pmatrix}
    \tilde{\wvec}\tilde{\wvec}^{\top} & u_0\tilde{\wvec}\\
    u_0\tilde{\wvec}^{\top} & u_0^2
\end{pmatrix}\right]^{-1}
\E\left[\begin{pmatrix}
    \tilde{\wvec}\tilde{\wvec}^{\top} \\ u_0\tilde{\wvec}^{\top}  
 \end{pmatrix} \right]}_{=:\Bvec_2} - \Bvec_2^{\top} +
 \underbrace{\E\left[\begin{pmatrix}
    \rvec\tilde{\wvec}^{\top} & u_0\rvec  
 \end{pmatrix} \right]
 \E\left[
 \begin{pmatrix}
    \tilde{\wvec}\tilde{\wvec}^{\top} & u_0\tilde{\wvec}\\
    u_0\tilde{\wvec}^{\top} & u_0^2
\end{pmatrix}\right]^{-1}
\E\left[\begin{pmatrix}
    \tilde{\wvec}\rvec^{\top} \\ u_0\rvec^{\top}  
 \end{pmatrix} \right]}_{=:\Bvec_3}\\
\stackrel{\text{(a)}}{=}&
\E\left[\tilde{\wvec}\tilde{\wvec}^{\top} \right]
\end{align*}
To show (a), we argue as follows.
First of all, 
$$
\E\left[
\rvec\tilde{\wvec}^{\top}\right]~=~
\E_{ z}\left[\rvec\left(\sum_{j\in \Kscr}
\E\left[\ind_{\{y = j\}} - g_{j}(\x, z; \thetavec^*)|\x, z
\right]\x_{(j)}\right)^{\top}
\right] = \Ovec
$$
and similarly $\E\left[
u_0\rvec\right] = 0$ so $\Bvec_2 = \Bvec_3 = \Ovec$.
Also, 
$$
\Bvec_1 ~=~ \begin{pmatrix}
  \Ivec & \Ovec   
\end{pmatrix}
\E\left[\begin{pmatrix}
    \tilde{\wvec}\tilde{\wvec}^{\top} \\ u_0\tilde{\wvec}^{\top}  
 \end{pmatrix} \right]
~=~\E\left[\tilde{\wvec}\tilde{\wvec}^{\top} \right].
$$
This proves (a).

Further, 
$$
\E\left[\wvec \wvec^{\top} \right] ~=~ \E\left[(\tilde{\wvec} - \rvec)(\tilde{\wvec} - \rvec)^{\top} \right] ~\stackrel{\text{(b)}}{=}~
\E\left[\tilde{\wvec}\tilde{\wvec}^{\top} \right] + \E\left[\rvec\rvec^{\top} \right],
$$
where (b) follows because $\E\left[
\rvec\tilde{\wvec}^{\top}\right] = \Ovec$.
This proves that 
$\check{\Jvec} - \Gammavec\Lambdavec\Gammavec^{\top} = \E\left[\rvec\rvec^{\top} \right]$. 
Also, 
using the fact that 
$$
\rvec~=~\sum_{j\in \Kscr}\big(\sigma_{j}(\x; \betavec^*) 
 - \Prob(y = j\,|\,\x)\big)\x_{(j)} + \big(\Prob(y = j\,|\,\x) - g_{j}(\x, z; \thetavec^*)\big)\x_{(j)}
$$
and that $\E_{z}[g_{j}(\x, z; \thetavec^*)|\x] = \Prob(y = j\,|\,\x),$ 
we have the decomposition. 
\hfill $\blacksquare$
\vspace{3mm}

\vspace{3mm}
